% Generated draft: requires mathematical and source review.
\documentclass[11pt]{article}
\usepackage[margin=25mm]{geometry}
\usepackage{fontspec}
\setmainfont{Latin Modern Roman}
\usepackage{amsmath,amssymb,mathtools,unicode-math}
\setmathfont{Latin Modern Math}
\usepackage{enumitem,needspace,xurl,hyperref}
\setlength{\parindent}{0pt}
\setlength{\parskip}{5pt}
\setlength{\emergencystretch}{3em}
\newcommand{\N}{\mathbb{N}}
\newcommand{\Z}{\mathbb{Z}}
\newcommand{\Q}{\mathbb{Q}}
\newcommand{\R}{\mathbb{R}}
\newcommand{\C}{\mathbb{C}}
\newcommand{\F}{\mathbb{F}}
\newcommand{\A}{\mathbb{A}}
\newcommand{\PP}{\mathbb{P}}
\newcommand{\E}{\mathbb{E}}
\newcommand{\eps}{\varepsilon}
\newcommand{\dd}{\,\mathrm{d}}
\begin{document}
\section*{UnsolvedMath ID: 30005748}
\textbf{Source problem code:} catalogue entry
\subsection{Definitions and mathematical statement}
The supplied record concerns implicit propositional proof systems, but does not define that class, the construction \([P,P]\), or the technical meaning of a jump in proof strength. Retain the source notation
\[
\operatorname{imp}(P):=[P,P].
\]
The intended conjecture, on the universal reading of the supplied statement, is: for every implicit propositional proof system \(P\) within the source's intended scope, \(\operatorname{imp}(P)\) is a jump in proof strength.

This is a schematic statement requiring clarification from the source. In particular, the admissible systems and any additional hypotheses on \(P\), the definition of \([P,P]\), the comparison relation on proof systems, and the conditions defining a jump are absent from the supplied excerpt. Consequently, no precise assertion about strict increase, simulation, or other properties of a jump can be substituted without additional evidence. The excerpt itself says only that \(\operatorname{imp}(P):=[P,P]\) is conjectured to be a jump; it does not explicitly state the domain or quantifier for \(P\).
\subsection{Short English statement}
The supplied conjecture asks whether the implicit construction \(\operatorname{imp}(P)=[P,P]\) gives a jump in proof strength for each admissible implicit propositional proof system \(P\). Its precise mathematical content remains unresolved because the supplied excerpt omits the construction, the admissibility hypotheses, and the definition of a jump.
\subsection{Sources}
\begin{itemize}
\item Supplied local UnsolvedMath registry record: permanent ID \(30005748\), source problem code \texttt{catalogue entry}, title ``Jump Property of Implication Proof Systems.''
\item Supplied report citation: \emph{Mathematical Logic: Proof Theory, Constructive Mathematics} (2024), Oberwolfach Reports, pp.\ 3013--3061, DOI \url{https://doi.org/10.4171/owr/2023/53}. The literature-source list labels this reference ``Oberwolfach Report 53/2023.'' The supplied excerpt associates the conjecture with Pavel Pudl\'ak's contribution ``Consistency, implicit proofs, and cut-elimination,'' p.\ 3023. The record gives proposed year 2023 and leaves the proposer field empty.
\item Supplied literature reference: Egidy, \emph{Recursive Jump Operators and Optimal Proof Systems}, \url{https://doi.org/10.4230/LIPIcs.ICALP.2026.88}.
\item Supplied literature assessment dated exactly \texttt{2026-08-22}: the record labels the problem open, states that the report presents \(\operatorname{imp}(P)=[P,P]\) as a conjecture, and says that a 2026 study still describes jump existence for \(\mathrm{TAUT}\) as open, with no resolution of the extracted conjecture available in that assessment. This is supplied dated curation, not independently verified here; the record itself requests independent verification of time-sensitive status and cited claims.
\item The registry marks the statement \texttt{corrected\_verified} and reports a check against the official Oberwolfach report. That verification claim is supplied metadata and has not been independently confirmed here.
\end{itemize}
\end{document}
